\documentclass[9pt,a4paper]{article}

% Basic packages
\usepackage[margin=1in]{geometry}
\usepackage{amsmath,amssymb,amsthm}
\usepackage{tikz}
\usetikzlibrary{arrows.meta,positioning,shapes}
\usepackage{xcolor}
\usepackage{hyperref}
\usepackage{microtype}
\usepackage{float}

% Unicode font support (required for Agda output - compile with lualatex)
\usepackage{fontspec}
\usepackage{unicode-math}

% Use DejaVu for excellent Unicode coverage (subscripts, arrows, etc.)
\setmainfont{DejaVu Serif}
\setsansfont{DejaVu Sans}
\setmonofont{DejaVu Sans Mono}[Scale=MatchLowercase]
\setmathfont{STIX Two Math}

% Agda literate support
\usepackage{agda}

% Override Agda font style: smaller size + DejaVu Sans for Unicode
% Must come AFTER loading agda.sty
\renewcommand{\AgdaFontStyle}[1]{{\small\sffamily #1}}
\renewcommand{\AgdaKeywordFontStyle}[1]{{\small\sffamily\bfseries #1}}
\renewcommand{\AgdaBoundFontStyle}[1]{{\small\sffamily\itshape #1}}
\renewcommand{\AgdaStringFontStyle}[1]{{\small\ttfamily #1}}
\renewcommand{\AgdaCommentFontStyle}[1]{{\small\ttfamily #1}}

% Custom color scheme: "Nord-inspired" - cool, muted tones
% More readable and modern than the default saturated colors
\definecolor{AgdaComment}      {HTML}{616E88}  % Muted blue-gray
\definecolor{AgdaPragma}       {HTML}{616E88}
\definecolor{AgdaKeyword}      {HTML}{5E81AC}  % Steel blue
\definecolor{AgdaString}       {HTML}{A3BE8C}  % Sage green
\definecolor{AgdaNumber}       {HTML}{B48EAD}  % Muted purple
\definecolor{AgdaSymbol}       {HTML}{4C566A}  % Dark gray

\definecolor{AgdaBound}                 {HTML}{2E3440}  % Near black
\definecolor{AgdaGeneralizable}         {HTML}{5E81AC}
\definecolor{AgdaInductiveConstructor}  {HTML}{8FBCBB}  % Teal
\definecolor{AgdaCoinductiveConstructor}{HTML}{88C0D0}  % Cyan
\definecolor{AgdaDatatype}              {HTML}{81A1C1}  % Light blue
\definecolor{AgdaField}                 {HTML}{B48EAD}  % Purple
\definecolor{AgdaFunction}              {HTML}{5E81AC}  % Steel blue
\definecolor{AgdaMacro}                 {HTML}{A3BE8C}  % Green
\definecolor{AgdaModule}                {HTML}{E0BA62}  % Gold
\definecolor{AgdaPostulate}             {HTML}{BF616A}  % Muted red
\definecolor{AgdaPrimitive}             {HTML}{81A1C1}  % Light blue
\definecolor{AgdaRecord}                {HTML}{81A1C1}  % Light blue
\definecolor{AgdaArgument}              {HTML}{4C566A}  % Dark gray

% Theorem environments
\theoremstyle{plain}
\newtheorem{theorem}{Theorem}
\newtheorem{lemma}[theorem]{Lemma}
\newtheorem{corollary}[theorem]{Corollary}
\theoremstyle{definition}
\newtheorem{definition}{Definition}
\theoremstyle{remark}
\newtheorem{remark}{Remark}

\title{\textbf{CSHRL Appendix} \\
Supplementary Proofs: Subsumption, Stream Isomorphism, and Stochastic Extension}

\author{Ricardo Corral-Corral}
\date{February 2026}

\begin{document}

\maketitle

\begin{abstract}
This appendix provides supplementary machine-verified proofs for the main CSHRL paper.
We establish three results:
(1)~the \textbf{Subsumption Theorem}---coinductive optimality implies classical argmax for any finite horizon;
(2)~the \textbf{Stream Isomorphism}---coinductive stream ordering is equivalent to pointwise dominance;
and (3)~the \textbf{Stochastic Subsumption Theorem}---extending argmax subsumption to stochastic MDPs via expected reward streams.
We conclude with a conjecture on first-order stochastic dominance (FOSD) for risk-sensitive extensions.
All proofs are literate Agda compiled with \texttt{--safe} and \texttt{--guardedness}.
\end{abstract}

%%%%%%%%%%%%%%%%%%%%%%%%%%%%%%%%%%%%%%%%%%%%%%%%%%%%%%%%%%%%%%%%%%%%%%%%%%%%%%%
\section{Introduction}
%%%%%%%%%%%%%%%%%%%%%%%%%%%%%%%%%%%%%%%%%%%%%%%%%%%%%%%%%%%%%%%%%%%%%%%%%%%%%%%

This document accompanies the main CSHRL paper with supplementary proofs.
After establishing common definitions (\S\ref{sec:prelim}), we present three machine-verified results and one conjecture:
\begin{enumerate}
\item \textbf{Subsumption Theorem} (\S\ref{sec:subsumption}): Coinductive optimality implies classical argmax optimality for any finite horizon $N$.
\item \textbf{Stream Isomorphism} (\S\ref{sec:iso}): $x \le_s y \Longleftrightarrow \forall n.\, x_n \le_r y_n$---the correspondence is bidirectional.
\item \textbf{Stochastic Subsumption} (\S\ref{sec:stoch-sub}): The subsumption theorem extends to stochastic MDPs via expected reward streams.
\item \textbf{FOSD Conjecture} (\S\ref{sec:fosd}): A proposed extension to first-order stochastic dominance for risk-sensitive RL.
\end{enumerate}

\begin{code}
{-# OPTIONS --safe --guardedness #-}

module CSHRL-Appendix where

open import Data.List using (List; map; foldr; []; _∷_)
open import Data.Nat using (ℕ; _⊔_)
open import Data.Product using (_×_; _,_; proj₁; proj₂)
open import Relation.Binary.PropositionalEquality using (_≡_; refl)
open import Codata.Guarded.Stream using (Stream; head; tail; tabulate)
\end{code}

%%%%%%%%%%%%%%%%%%%%%%%%%%%%%%%%%%%%%%%%%%%%%%%%%%%%%%%%%%%%%%%%%%%%%%%%%%%%%%%
\section{Preliminaries}
\label{sec:prelim}
%%%%%%%%%%%%%%%%%%%%%%%%%%%%%%%%%%%%%%%%%%%%%%%%%%%%%%%%%%%%%%%%%%%%%%%%%%%%%%%

We begin by recalling the core definitions from the main development: reward streams, the coinductive stream ordering $\le_s$, and the coinductive symmetric homomorphism record \AgdaRecord{CoindHomo}.

\begin{code}
module PreservationEquivalence
  (State Action Reward : Set)
  (step                : State → Action → State × Reward)
  (_≤ᵣ_                : Reward → Reward → Set)
  (max                 : Reward → Reward → Reward)
  (bottom              : Reward)
  (all-actions         : List Action)
  where

  -- Stream of rewards
  StreamR : Set
  StreamR = Stream Reward

  -- Optimal value at depth n
  max-list : List Reward → Reward
  max-list = foldr max bottom

  solve : State → ℕ → Reward
  solve s ℕ.zero    = max-list (map (λ a → proj₂ (step s a)) all-actions)
  solve s (ℕ.suc n) = max-list (map (λ a → solve (proj₁ (step s a)) n) all-actions)

  value : State → StreamR
  value s = tabulate (solve s)

  action-value : State → Action → StreamR
  head (action-value s a) = proj₂ (step s a)
  tail (action-value s a) = value (proj₁ (step s a))

  -- Coinductive stream ordering
  record _≤ₛ_ (x y : StreamR) : Set where
    coinductive
    field
      head≤ : head x ≤ᵣ head y
      tail≤ : tail x ≤ₛ tail y

  open _≤ₛ_ public

  -- The coinductive symmetric homomorphism
  record CoindHomo : Set₁ where
    field
      _≤ₐ_      : State → Action → Action → Set
      preserves : ∀ a b s → _≤ₐ_ s a b →
                  action-value s a ≤ₛ action-value s b
\end{code}

\begin{code}[hide]
  -- Auxiliary definitions (product formulation and η-equivalence)
  -- used only for module completeness; not discussed in the appendix.
  preserves-× : (State → Action → Action → Set) → Set
  preserves-× _≤ₐ_ = ∀ a b s → _≤ₐ_ s a b →
                     let v₁ = action-value s a
                         v₂ = action-value s b
                     in (head v₁ ≤ᵣ head v₂) × (tail v₁ ≤ₛ tail v₂)

  optimality-η : {_≤ₐ_ : State → Action → Action → Set} →
                 preserves-× _≤ₐ_ →
                 ∀ s a b → _≤ₐ_ s a b →
                 action-value s a ≤ₛ action-value s b
  head≤ (optimality-η pres s a b p) = proj₁ (pres a b s p)
  tail≤ (optimality-η pres s a b p) = proj₂ (pres a b s p)

  all-three-equal : (h : CoindHomo) →
                    ∀ s a b (p : CoindHomo._≤ₐ_ h s a b) →
                    let pres-× : preserves-× (CoindHomo._≤ₐ_ h)
                        pres-× a' b' s' r' = let q = CoindHomo.preserves h a' b' s' r'
                                             in head≤ q , tail≤ q
                        via-η = optimality-η pres-× s a b p
                        original = CoindHomo.preserves h a b s p
                    in (head≤ via-η ≡ head≤ original)
                     × (tail≤ via-η ≡ tail≤ original)
  all-three-equal _ _ _ _ _ = refl , refl
\end{code}

%%%%%%%%%%%%%%%%%%%%%%%%%%%%%%%%%%%%%%%%%%%%%%%%%%%%%%%%%%%%%%%%%%%%%%%%%%%%%%%
\section{Subsumption of Classical Argmax}
\label{sec:subsumption}
%%%%%%%%%%%%%%%%%%%%%%%%%%%%%%%%%%%%%%%%%%%%%%%%%%%%%%%%%%%%%%%%%%%%%%%%%%%%%%%

The main paper establishes that CSHRL's coinductive optimality \emph{subsumes} classical RL's argmax criterion. Here we provide the complete proof.

\subsection{The Bridge: Pointwise Dominance}

The key insight is that coinductive stream ordering ($\le_s$) implies \emph{pointwise} dominance at every index:

\begin{code}
  -- Extract the n-th element of a stream
  iter-head : ℕ → StreamR → Reward
  iter-head ℕ.zero    s = head s
  iter-head (ℕ.suc n) s = iter-head n (tail s)

  -- Pointwise dominance: at every index
  PointwiseDominance : StreamR → StreamR → Set
  PointwiseDominance x y = ∀ n → iter-head n x ≤ᵣ iter-head n y

  -- The bridge lemma
  ≤ₛ-to-pointwise : ∀ {x y} → x ≤ₛ y → PointwiseDominance x y
  ≤ₛ-to-pointwise p ℕ.zero    = head≤ p
  ≤ₛ-to-pointwise p (ℕ.suc n) = ≤ₛ-to-pointwise (tail≤ p) n
\end{code}

This lemma is the crucial connection: it converts the \emph{infinite} coinductive structure into a property we can use for \emph{finite} sums.

\subsection{The Full Subsumption Module}

The complete proof requires extending Core with reward arithmetic. The standalone implementation is in \texttt{src/appendix/Arithmetic.agda}. Here we show the module signature and key theorem:

\begin{code}
module SubsumptionTheorem
  (State Action Reward : Set)
  (step    : State → Action → State × Reward)
  (_≤ᵣ_    : Reward → Reward → Set)
  (max     : Reward → Reward → Reward)
  (bottom  : Reward)
  (actions : List Action)
  -- Arithmetic extension
  (_+ᵣ_      : Reward → Reward → Reward)
  (≤ᵣ-refl   : ∀ {r} → r ≤ᵣ r)
  (+ᵣ-mono-≤ : ∀ {a b c d} → a ≤ᵣ b → c ≤ᵣ d → (a +ᵣ c) ≤ᵣ (b +ᵣ d))
  where

  open PreservationEquivalence State Action Reward step _≤ᵣ_ max bottom actions

  -- Partial sum definition
  partial-sum : ℕ → StreamR → Reward
  partial-sum ℕ.zero    _ = bottom
  partial-sum (ℕ.suc n) s = head s +ᵣ partial-sum n (tail s)

  -- The subsumption theorem type
  SubsumesPartialSum : CoindHomo → Set
  SubsumesPartialSum homo = ∀ s a b N →
    CoindHomo._≤ₐ_ homo s a b →
    partial-sum N (action-value s a) ≤ᵣ partial-sum N (action-value s b)
\end{code}

\subsection{Proof Structure}

The proof proceeds in three steps:

\begin{enumerate}
\item \textbf{Preservation:} Given $a \le_a b$ in the ranking, \texttt{CoindHomo.preserves} yields \texttt{action-value s a $\le_s$ action-value s b}.

\item \textbf{Pointwise conversion:} Apply \texttt{≤ₛ-to-pointwise} to get $\forall n.\, r^a_n \le_r r^b_n$.

\item \textbf{Induction on horizon:} Prove \texttt{partial-sum-mono} by induction on $N$:
  \begin{itemize}
  \item Base ($N = 0$): Both sums equal \texttt{bottom}; use \texttt{≤ᵣ-refl}.
  \item Step ($N = k+1$): The sum is $r_0 + \sum_{t=1}^{k} r_t$. Apply \texttt{+ᵣ-mono-≤} to the head inequality and the inductive hypothesis on tails.
  \end{itemize}
\end{enumerate}

\begin{code}
  -- Helper for shifting pointwise to tails
  pointwise-tail : ∀ {x y} → PointwiseDominance x y → 
                   PointwiseDominance (tail x) (tail y)
  pointwise-tail pw n = pw (ℕ.suc n)

  -- Partial sum monotonicity (the inductive core)
  partial-sum-mono : ∀ N x y → PointwiseDominance x y → 
                     partial-sum N x ≤ᵣ partial-sum N y
  partial-sum-mono ℕ.zero _ _ _ = ≤ᵣ-refl
  partial-sum-mono (ℕ.suc n) x y pw = 
    +ᵣ-mono-≤ (pw ℕ.zero) 
              (partial-sum-mono n (tail x) (tail y) (pointwise-tail pw))

  -- The complete proof
  subsumes-partial-sum : (h : CoindHomo) → SubsumesPartialSum h
  subsumes-partial-sum h s a b N p = 
    partial-sum-mono N (action-value s a) (action-value s b) 
                     (≤ₛ-to-pointwise (CoindHomo.preserves h a b s p))
\end{code}

\subsection{Corollary: Argmax Identification}

An immediate corollary: if action $b$ is ranked highest ($\forall a.\, a \le_a b$), then $b$ has the highest partial sum for \emph{any} horizon:

\begin{code}
  argmax-subsumed : (h : CoindHomo) → ∀ s b →
    (∀ a → CoindHomo._≤ₐ_ h s a b) →
    ∀ a N → partial-sum N (action-value s a) ≤ᵣ partial-sum N (action-value s b)
  argmax-subsumed h s b b-top a N = subsumes-partial-sum h s a b N (b-top a)
\end{code}

This is the formal statement that \textbf{CSHRL rankings identify the classical argmax}---not by computing expected returns, but as a structural consequence of the coinductive homomorphism property.

%%%%%%%%%%%%%%%%%%%%%%%%%%%%%%%%%%%%%%%%%%%%%%%%%%%%%%%%%%%%%%%%%%%%%%%%%%%%%%%
\section{Stream Isomorphism}
\label{sec:iso}
%%%%%%%%%%%%%%%%%%%%%%%%%%%%%%%%%%%%%%%%%%%%%%%%%%%%%%%%%%%%%%%%%%%%%%%%%%%%%%%

The Subsumption Theorem uses one direction of a deeper correspondence: coinductive stream ordering implies pointwise dominance (\AgdaFunction{≤ₛ-to-pointwise}).  We now prove the \emph{converse}---pointwise dominance implies coinductive ordering---establishing a full isomorphism.

\subsection{The Converse and the Isomorphism}

\begin{code}
  -- Converse: pointwise dominance implies coinductive ordering
  pointwise-to-≤ₛ : ∀ {x y} → PointwiseDominance x y → x ≤ₛ y
  head≤ (pointwise-to-≤ₛ pw) = pw ℕ.zero
  tail≤ (pointwise-to-≤ₛ pw) = pointwise-to-≤ₛ (pointwise-tail pw)

  _⇔_ : Set → Set → Set
  A ⇔ B = (A → B) × (B → A)

  -- The isomorphism: combining both directions
  stream-iso : ∀ x y → (x ≤ₛ y) ⇔ (PointwiseDominance x y)
  stream-iso x y = ≤ₛ-to-pointwise , pointwise-to-≤ₛ
\end{code}

\noindent This establishes:
\[
x \le_s y \;\Longleftrightarrow\; \forall n.\, x_n \le_r y_n
\]

The proof of \AgdaFunction{pointwise-to-≤ₛ} is corecursive: to construct $x \le_s y$ from pointwise dominance, we extract the head inequality from \texttt{pw\;0} and corecursively prove the tail inequality by shifting the pointwise witness via \AgdaFunction{pointwise-tail}.

\subsection{Implications}

The isomorphism has two consequences for the Finder algorithm, which defines rankings via trace comparison at finite depth:
\begin{enumerate}
\item \textbf{Finder $\to$ CoindHomo:} If traces dominate at all depths, streams dominate coinductively.
\item \textbf{CoindHomo $\to$ Finder:} If streams dominate coinductively, traces dominate at all depths.
\end{enumerate}

\noindent Together, these show that the Finder's ranking (at sufficient depth) \emph{exactly} characterizes the coinductive property, not merely an approximation.  The ``symmetric'' in CSHRL is thus formally justified: the correspondence between rankings and stream dominance is an isomorphism, not just a one-way homomorphism.

The standalone module \texttt{src/appendix/StreamIsomorphism.agda} provides the same proof in a reusable parameterized form.

%%%%%%%%%%%%%%%%%%%%%%%%%%%%%%%%%%%%%%%%%%%%%%%%%%%%%%%%%%%%%%%%%%%%%%%%%%%%%%%
\section{Stochastic Subsumption Theorem}
\label{sec:stoch-sub}
%%%%%%%%%%%%%%%%%%%%%%%%%%%%%%%%%%%%%%%%%%%%%%%%%%%%%%%%%%%%%%%%%%%%%%%%%%%%%%%

The deterministic subsumption theorem (\S\ref{sec:subsumption}) shows that coinductive optimality implies classical argmax for any finite horizon.  A natural question: does this extend to \emph{stochastic} MDPs?

\subsection{The Key Insight: Linearity of Expectation}

Classical RL maximizes \emph{expected} cumulative reward:
\[
\text{argmax}_a \; \mathbb{E}\left[\sum_{t=0}^{N-1} r_t(a)\right]
\]

By linearity of expectation:
\[
\mathbb{E}\left[\sum_{t=0}^{N-1} r_t\right] = \sum_{t=0}^{N-1} \mathbb{E}[r_t]
\]

This means comparing partial sums of \emph{expected} rewards is equivalent to comparing \emph{expected} partial sums---exactly what classical RL optimizes. The deterministic subsumption proof therefore applies directly to expected reward streams.

\subsection{Stochastic Pointwise Dominance}

For subsumption, we use \textbf{pointwise expected dominance} (stronger than lexicographic):

\begin{code}[hide]
module StochasticSubsumptionSection
  (State Action Reward : Set)
  (step : State → Action → List (Reward × ℕ))
  (_≤ᵣ_ : Reward → Reward → Set)
  (≤ᵣ-refl : ∀ {r} → r ≤ᵣ r)
  (_+ᵣ_ : Reward → Reward → Reward)
  (_*ᵣ_ : ℕ → Reward → Reward)
  (zeroᵣ : Reward)
  (+ᵣ-mono-≤ : ∀ {a b c d} → a ≤ᵣ b → c ≤ᵣ d → (a +ᵣ c) ≤ᵣ (b +ᵣ d))
  where

  open import Codata.Guarded.Stream using (Stream; head; tail)
  open import Relation.Binary.PropositionalEquality using (_≡_)
  StreamR : Set
  StreamR = Stream Reward
\end{code}

\begin{code}
  -- Pointwise expected stream dominance (for subsumption)
  record _≤ₛ-expected_ (x y : StreamR) : Set where
    coinductive
    field
      head≤ : head x ≤ᵣ head y
      tail≤ : tail x ≤ₛ-expected tail y
\end{code}

Compare with the \emph{lexicographic} version from the main stochastic theory:

\begin{code}
  -- Lexicographic: tail comparison is CONDITIONAL on head equality
  record _≤ₛ-lex_ (x y : StreamR) : Set where
    coinductive
    field
      head≤ : head x ≤ᵣ head y
      tail≤ : head x ≡ head y → tail x ≤ₛ-lex tail y
\end{code}

The pointwise version requires tail dominance \emph{unconditionally}, while lexicographic only requires it when heads are equal. Pointwise is strictly stronger: it implies lexicographic, but not vice versa.

\subsection{The Stochastic Subsumption Theorem}

\begin{code}[hide]
  open _≤ₛ-expected_ public

  iter-head : ℕ → StreamR → Reward
  iter-head ℕ.zero    s = head s
  iter-head (ℕ.suc n) s = iter-head n (tail s)

  PointwiseDominance : StreamR → StreamR → Set
  PointwiseDominance x y = ∀ n → iter-head n x ≤ᵣ iter-head n y

  ≤ₛ-expected-to-pointwise : ∀ {x y} → x ≤ₛ-expected y → PointwiseDominance x y
  ≤ₛ-expected-to-pointwise p ℕ.zero    = head≤ p
  ≤ₛ-expected-to-pointwise p (ℕ.suc n) = ≤ₛ-expected-to-pointwise (tail≤ p) n

  partial-sum : ℕ → StreamR → Reward
  partial-sum ℕ.zero    _ = zeroᵣ
  partial-sum (ℕ.suc n) s = head s +ᵣ partial-sum n (tail s)

  pointwise-tail : ∀ {x y} → PointwiseDominance x y → PointwiseDominance (tail x) (tail y)
  pointwise-tail pw n = pw (ℕ.suc n)
\end{code}

The proof structure mirrors the deterministic case exactly:

\begin{code}
  -- Partial sum respects pointwise dominance
  partial-sum-mono : ∀ N x y → PointwiseDominance x y →
                     partial-sum N x ≤ᵣ partial-sum N y
  partial-sum-mono ℕ.zero    _ _ _  = ≤ᵣ-refl
  partial-sum-mono (ℕ.suc n) x y pw =
    +ᵣ-mono-≤ (pw ℕ.zero)
              (partial-sum-mono n (tail x) (tail y) (pointwise-tail pw))
\end{code}

\begin{theorem}[Stochastic Subsumption]
If action rankings satisfy pointwise expected stream dominance, then for any finite horizon $N$, the ranked-higher action has higher expected cumulative reward:
\[
a \le_{\text{rank}} b \;\wedge\; \mathbb{E}[v(a)] \le_s^{\text{pw}} \mathbb{E}[v(b)] \;\Longrightarrow\; \sum_{t=0}^{N-1} \mathbb{E}[r_t(a)] \le \sum_{t=0}^{N-1} \mathbb{E}[r_t(b)]
\]
\end{theorem}

The standalone implementation is in \texttt{src/appendix/StochasticSubsumption.agda}.

\subsection{Two Orderings, Two Purposes}

\begin{center}
\begin{tabular}{|l|l|l|}
\hline
\textbf{Ordering} & \textbf{Purpose} & \textbf{Role} \\
\hline
Lexicographic ($\le_s^{\text{lex}}$) & Trace-based learning, ranking semantics & \textbf{Primary} \\
Pointwise ($\le_s^{\text{pw}}$) & Argmax subsumption, classical RL connection & Auxiliary \\
\hline
\end{tabular}
\end{center}

\begin{itemize}
\item \textbf{Lexicographic (primary):} The natural choice for CSHRL---aligned with finite trace learning and captures the ``first difference decides'' semantics. This is the ordering used in \texttt{StochasticCoindHomo}.
\item \textbf{Pointwise (auxiliary):} A more restrictive ordering used \emph{only} to establish the subsumption connection to classical RL. Many well-behaved MDPs satisfy pointwise when the ranking is sound.
\end{itemize}

The relationship: $\text{Pointwise} \Rightarrow \text{Lexicographic}$, but not vice versa. Pointwise is more restrictive (requires dominance at \emph{all} timesteps), while lexicographic is more permissive (only requires comparison when heads are equal). The lexicographic choice is intentional: it matches how finite trace comparison naturally extends to the infinite case.

%%%%%%%%%%%%%%%%%%%%%%%%%%%%%%%%%%%%%%%%%%%%%%%%%%%%%%%%%%%%%%%%%%%%%%%%%%%%%%%
\section{Conjecture: First-Order Stochastic Dominance Isomorphism}
\label{sec:fosd}
%%%%%%%%%%%%%%%%%%%%%%%%%%%%%%%%%%%%%%%%%%%%%%%%%%%%%%%%%%%%%%%%%%%%%%%%%%%%%%%

\subsection{Beyond Expected Values: Risk-Sensitive Orderings}

The stochastic subsumption theorem (\S\ref{sec:stoch-sub}) compares \emph{expected values}---collapsing distributions to their means. A richer question: can we compare \emph{full distributions} over reward streams?

\textbf{First-Order Stochastic Dominance (FOSD)} is a classical ordering on probability distributions: distribution $\mu$ FOSD-dominates $\nu$ (written $\mu \ge_{\text{FOSD}} \nu$) iff for all thresholds $r$:
\[
P_\mu(X \le r) \le P_\nu(X \le r)
\]
Equivalently, the CDF of $\mu$ lies everywhere below the CDF of $\nu$. FOSD is strictly stronger than expected-value comparison: $\mu \ge_{\text{FOSD}} \nu \Rightarrow \mathbb{E}[\mu] \ge \mathbb{E}[\nu]$, but not vice versa.

\subsection{The Conjecture}

Let $\mu, \nu$ be probability distributions over infinite reward streams. Define:
\begin{itemize}
\item \textbf{Coinductive FOSD:} $\mu \le_s \nu$ iff the marginal distribution of heads satisfies FOSD, and the conditional tail distributions satisfy $\mu_{\text{tail}} \le_s \nu_{\text{tail}}$ (corecursively).
\item \textbf{Pointwise FOSD:} $\mu \le_{\text{pw}} \nu$ iff for all timesteps $n$ and thresholds $r$, $P_\mu(X_n \le r) \ge P_\nu(X_n \le r)$.
\end{itemize}

\begin{center}
\textbf{Conjecture (Stochastic Isomorphism):} $\mu \le_s \nu \;\Longleftrightarrow\; \mu \le_{\text{pw}} \nu$
\end{center}

If true, this would extend CSHRL to stochastic environments with a \emph{risk-sensitive} interpretation: a policy that stochastically dominates is preferred by \emph{all} risk-averse decision makers, not just those maximizing expected return.

\subsection{Current Infrastructure and Remaining Gap}

The Giry monad foundation is now in place (\texttt{CSHRL.Probability.Finite}): finite distributions as weighted lists, with monadic operations (\texttt{pure}, \texttt{>>=}, \texttt{fmap}) and expected value computation. The stochastic extension (\texttt{CSHRL.Core.Stochastic}) uses this to define expected reward streams and lexicographic comparison.

However, proving the FOSD isomorphism conjecture requires additional infrastructure not yet implemented:
\begin{itemize}
\item \textbf{CDF computation:} Cumulative distribution functions for \texttt{Dist Reward}, requiring decidable ordering on rewards.
\item \textbf{FOSD comparison:} $d_1 \le_{\text{FOSD}} d_2$ iff $\forall r.\, \text{cdf}(d_1, r) \ge \text{cdf}(d_2, r)$.
\item \textbf{Conditional tail distributions:} Preserving correlation structure through coinductive unfolding.
\end{itemize}

The current expected-value approach (\S\ref{sec:stoch-sub}) is sufficient for classical RL subsumption.  The FOSD extension would enable \emph{risk-sensitive} comparisons---a distinct research direction beyond expected-value optimization.

\end{document}


